\subsection*{Algorithm Track Benchmark Harness}
    The NeuroBench algorithm benchmarks are wrapped in a harness which standardizes the benchmark interfaces. The harness provides benchmark users with a consistent framework for loading data, processing data and model outputs, and calculating and reporting metrics, thereby ensuring fair and standard comparisons of the results. It is built with straightforward interfaces which are designed to be extended with new frameworks, algorithms, and tasks. The benchmark harness is open-source for use and development (\url{https://github.com/NeuroBench/neurobench}).
    %\footnote{\url{https://github.com/NeuroBench/neurobench}}
    
    The components of the algorithm benchmark harness are summarized in Figure~\ref{fig:AlgTrack_SA}.
    % Note that the execution flow currently assumes a digital time-stepped execution of the algorithm, but can be extended to continuous values such as in analog design in the future
    \textit{Datasets} are loaded in a common format and pass through \textit{Processors} to be pre-processed. The \textit{Model} generates predictions based on the processed data, and \textit{Accumulators} post-process the predictions, for instance to accumulate spikes and transform to labels. Static metrics of algorithm footprint and connection sparsity are calculated via model analysis, while metrics of correctness, activation sparsity, and synaptic operations are calculated using predictions and model execution traces. For benchmark users, task evaluation simply involves utilizing the existing dataloaders, processors, and metrics within the harness and wrapping their own code to fit the standard interfaces.
 
    Currently, the harness and all baseline models are built using PyTorch~\cite{paszke19pytorch} or frameworks based on it, such as snnTorch~\cite{eshraghian2021training} and SpikingJelly~\cite{SpikingJelly}. Due to its modular structure and simple interfaces, the harness can grow to be compatible with further neuromorphic tools such as Lava~\cite{lava} and Fugu~\cite{aimone2019composing}.
    Furthermore, it also supports the extension of data and metric pipelines in order to implement additional benchmark tasks. \revision{Widely validated benchmarks in keyword~\cite{warden2018speech} and gesture classification~\cite{Amir_etal17_lowpowe}, which are foundational in neuromorphic and conventional research~\cite{blouw2019, Cramer_2022, fang2021incorporating, massa2021efficient}, have been incorporated into the harness to complement the novel tasks in the NeuroBench v1.0 suite.} Any novel or existing benchmarks can make use of the harness infrastructure for open reproducibility, and also to garner interest in the community towards long-term task support and appearance in NeuroBench-affiliated leaderboards and challenge events.
    % As NeuroBench is a community-driven project, we welcome interested developers to contribute to the benchmark harness through the associated code repository. 

\subsubsection*{Algorithm Track Limitations and Further Extensions}

Before diving into the baseline results, it is worth discussing several possible improvements to the NeuroBench algorithm track framework in its current form. Specifically, the initial iteration of metrics is restricted to the assumption of digital, time-stepped algorithm execution. While complexity analysis of such prototypes can serve as an intermediate step for solutions intended for analog or continuous time deployment, the metric measurements are not yet defined for those execution settings. Informed by further benchmark implementations, future versions of NeuroBench will extend inclusiveness by expanding measurement protocols to include such algorithms.

\revision{Furthermore, the synaptic operations metric, intended to capture model computation cost, currently does not account for neuron updates. The dynamics of neuron models, including mechanisms like leakage and reset, can vary heavily in complexity. However, counting the number and type of operations from neuron updates, as well as estimating their overall costs, depends on the specific arithmetic or circuit implementation. Thus, they are not accounted for in the broader algorithmic complexity metrics. The algorithmic metric framework can be extended with solution-specific metrics that assume a particular implementation platform to estimate neuron update costs, which have been previously defined~\cite{Lemaire_2023}. These estimates can then be combined with the total number of neuron updates per model computation to measure overall network operation complexity during evaluation.}

% Solution-specific metrics that assume a particular implementation platform, as have been defined previously~\cite{Lemaire_2023}, can be used to estimate neuron update costs. These estimates can then be combined with the total number of neuron updates per model computation to measure overall neuron operation complexity during evaluation.

Data pre- and post-processing can also amount to significant costs not yet captured in the NeuroBench algorithm track metrics. Such costs are, however, captured in the deployed metrics of the system track, which accounts for data processing hardware as part of the overall system during performance and efficiency measurements. Data processing metrics will be added as a separate complexity category for the algorithm track benchmark in the future.

The v1.0 algorithm track benchmark suite is also intended to expand in the future. This could include covering further data modalities such as inertial measurement unit (IMU) sensing~\cite{fra22har} and extending to closed-loop sensorimotor tasks to demonstrate embodied intelligence. As with the initial benchmarks, further tasks will undergo approval and development by the open NeuroBench community before being included in a future versioned benchmark suite.

% Several improvements are possible on the benchmark metrics, tasks and harness. The first-iteration algorithm track metrics are currently restricted to the assumption of a digital time-stepped execution of the algorithm. Complexity analysis of digital time-stepped prototypes can be used as an intermediate step for solutions intended for analog or continuous time deployment, but the metric measurements are not yet defined for such execution. Informed by further benchmark implementations, future versions of NeuroBench will extend inclusiveness by expanding measurement protocols to include such algorithms.

% Furthermore, the synaptic operations metric, intended to capture the model computation cost, currently does not account for neuron updates. The dynamics of neuron models, including mechanisms such as leakage and reset, can vary heavily in complexity. However, the number and type of operations of neuron updates, as well as their overall costs, depend on arithmetic or circuit implementation, and thus are not accounted for in the broader algorithmic complexity metrics. Solution-specific metrics which assume a particular implementation platform, which have been defined previously~\cite{Lemaire_2023}, can be used to estimate neuron update costs. These can be combined with the total number of neuron updates per model computation in order to measure neuron operation complexity during evaluation.

% Data pre- and post-processing can amount to be significant. These costs are also not yet captured in the NeuroBench algorithm track metrics. Such costs are, however, captured in the deployed metrics of the system track, which accounts for data processing hardware as part of the overall system during performance and efficiency measurements. Data processing metrics will be added as a separate complexity category for the algorithm track in the future.

% The v1.0 algorithm track benchmark suite is also intended to be expanded in the future, for instance covering further data modalities such as inertial measurement unit (IMU) sensing~\cite{fra22har} and extending to closed-loop sensorimotor tasks to demonstrate embodied intelligence. As with the initial benchmarks, further tasks will undergo approval and development by the open NeuroBench community to be included in a future versioned benchmark suite.